\documentclass[11pt,a4paper,twocolumn]{article}

% ---- Packages ------------------------------------------------------
\usepackage[utf8]{inputenc}
\usepackage[T1]{fontenc}
\usepackage{lmodern}
\usepackage[margin=1in]{geometry}
\usepackage{microtype}
\usepackage{authblk}
\usepackage{abstract}
\usepackage{amsmath,amssymb}
\usepackage{booktabs}
\usepackage{tabularx}
\usepackage{enumitem}
\usepackage{xcolor}
\usepackage{listings}
\usepackage{verbatim}
\usepackage{hyperref}
\usepackage{cleveref}
\usepackage{fancyhdr}

\hypersetup{
    colorlinks=true,
    linkcolor=teal,
    citecolor=teal,
    urlcolor=teal,
    pdftitle={Vouch Protocol: Cryptographic Identity and Continuous State Verifiability for Autonomous AI Agents},
    pdfauthor={Ramprasad Anandam Gaddam}
}

% Allow LaTeX a small emergency stretch to absorb unbreakable tokens
% (long \verb|...| and \url{...} blocks) without overfull-hbox overflow.
\emergencystretch=3em
\tolerance=2000

% Make \url{...} break aggressively at slashes / hyphens / dots so long
% repository URLs do not run past the right margin.
\makeatletter
\g@addto@macro{\UrlBreaks}{\UrlOrds}
\makeatother

% ---- Listings styling ----------------------------------------------
\definecolor{codebg}{HTML}{F5F5F0}
\definecolor{codekw}{HTML}{1F4E79}
\definecolor{codestr}{HTML}{7C2D3A}
\definecolor{codecom}{HTML}{6B6B78}

\lstset{
    basicstyle=\ttfamily\footnotesize,
    backgroundcolor=\color{codebg},
    keywordstyle=\color{codekw}\bfseries,
    stringstyle=\color{codestr},
    commentstyle=\color{codecom}\itshape,
    breaklines=true,
    columns=fullflexible,
    showstringspaces=false,
    frame=single,
    framerule=0pt,
    framexleftmargin=4pt,
    framexrightmargin=4pt,
    framextopmargin=3pt,
    framexbottommargin=3pt,
    aboveskip=0.6em,
    belowskip=0.6em,
    captionpos=b
}
\lstdefinelanguage{json}{
    string=[s]{"}{"},
    stringstyle=\color{codestr},
    comment=[l]{//},
    commentstyle=\color{codecom}\itshape,
    showstringspaces=false
}

% ---- Header / Footer -----------------------------------------------
\pagestyle{fancy}
\fancyhf{}
\fancyhead[L]{\small Vouch Protocol}
\fancyhead[R]{\small Gaddam}
\fancyfoot[C]{\thepage}
\renewcommand{\headrulewidth}{0.4pt}

% ---- Title / Author ------------------------------------------------
\title{
    Vouch Protocol: Cryptographic Identity and Continuous State Verifiability for Autonomous AI Agents
}
\author[1]{Ramprasad Anandam Gaddam}
\affil[1]{Independent (open-source maintainer, Vouch Protocol). \\
\texttt{ram@vouch-protocol.com} \\
\url{https://github.com/vouch-protocol/vouch}}
\date{Version v0.1 \quad May 2026 \\
\emph{Licensed under CC BY 4.0} \\[4pt]
\small Vouch-signed by Vouch Protocol: \url{https://vch.sh/arxiv-1}}

% ====================================================================
\begin{document}
\maketitle

% ---- Abstract ------------------------------------------------------
\begin{abstract}
Autonomous AI agents now make decisions that matter: they move money, submit regulated filings, read clinical records, commit code. The credentials those agents use to do this (API keys, OAuth bearer tokens, shared service accounts) were designed for a human clicking through a session. They prove access. They do not prove who authorized that action, what the agent intends to do, or whether the agent is still trustworthy by the time it acts.

We present the \textbf{Vouch Protocol}: an open specification and a reference implementation in Python, TypeScript, and Go for cryptographically identifying autonomous AI agents and binding every action they take to a signed Verifiable Credential. Each credential carries the agent's identity, the specific intent (action, target, resource), the delegation chain from the original human principal down to the agent, and a freshness window after which the credential stops being valid.

Per-action signing is necessary but not sufficient. Once an agent has been admitted, the next question is whether it is still behaving correctly \emph{right now}, after we let it through the door. The protocol's \emph{State Verifiability layer} answers that with four mechanisms: (1) trust entropy decay, where credential trust falls exponentially over time and must be renewed before consequential actions; (2) behavioral attestation digests, lightweight per-interval summaries of the agent's API calls, resources accessed, and intent drift; (3) canary commit/reveal chains, providing cryptographic detection of silent agent failure or substitution; and (4) $M$-of-$N$ validator quorum, distributing trust evaluation across role-specialized validators (policy, behavior, budget).

The protocol is built on W3C Verifiable Credentials Data Model 2.0 with Data Integrity proofs: the \texttt{eddsa-jcs-2022} cryptosuite, which signs Ed25519 over RFC 8785 JCS-canonicalized payloads. For the post-quantum transition, the protocol defines a \emph{dual-proof profile} in which a credential carries two independent Data Integrity proofs on the same JCS-canonicalized bytes: one \texttt{eddsa-jcs-2022} proof (classical Ed25519) and one \texttt{mldsa44-jcs-2026} proof (ML-DSA-44, FIPS 204). Verifiers can downgrade gracefully, and the wire format requires no Vouch-specific composite cryptosuite identifier.

This paper covers the credential format, the Identity Sidecar pattern that isolates an agent's signing key from the LLM process, delegation-chain construction with resource narrowing and depth-limited validation, the BitstringStatusList revocation mechanism, the Heartbeat Protocol's renewal cycle, the dual-proof post-quantum profile and its same-canonical-bytes property, and the State Verifiability runtime. We then work through the adversarial cases that motivated the design (prompt injection, key exfiltration, replay across resources, post-quantum cryptographic obsolescence) and what the protocol does about each. Cross-language test vectors show that three independent implementations produce byte-identical credentials.

The implementation is open-source under Apache 2.0. Sixty defensive prior-art disclosures accompany the specification under CC0 to keep design innovations openly available.

\smallskip
\noindent\textbf{Index Terms:} AI agent identity, Verifiable Credentials, decentralized identifiers, post-quantum cryptography, dual-proof signatures, JCS canonicalization, prompt injection, delegation, continuous attestation, behavioral provenance.
\end{abstract}

% ====================================================================
\section{Introduction}

\subsection{The Accountability Gap}

Modern AI agents (LangChain pipelines, CrewAI swarms, AutoGPT-style autonomous workers, Model Context Protocol-driven assistants) increasingly take actions that human-level credentials were never designed to authorize. An LLM-driven agent submitting an insurance claim, placing a stock order, querying a patient's electronic health record, or auto-committing code to a production branch is performing a real action with real consequences. When that action requires audit, dispute resolution, or regulatory review, the only artefact available is a log entry naming an API key. The key is shared, was rotated three months ago, and gives no insight into which specific agent generation, which orchestration layer, which prompt, or which user-level principal authorized the action.

The accountability gap has three structural sources:

\begin{enumerate}[label=\textbf{\arabic*.},leftmargin=*]
\item \textbf{Identity opacity.} Bearer credentials (API keys, JWTs, OAuth access tokens) prove access by possession, not by signature over an intent. An entity in possession of the bearer can take any action the bearer permits; nothing in the credential reflects which specific entity, in which specific context, with which specific upstream authorization, did so.

\item \textbf{Intent omission.} Even when an action is logged, the intent it represented is recorded in application-specific log formats with no cryptographic binding. A log entry ``\texttt{agent-prod-2 called submit\_claim}'' can be forged by anyone with access to the log writer. The action's parameters, the resource it targeted, and the authorization chain are not cryptographically bound to the agent's identity.

\item \textbf{Trust monotonicity.} Existing identity systems trust credentials until they are explicitly revoked. A compromised agent continues operating until a human detects misbehavior and pushes a revocation. For autonomous agents, which may operate in regulated, latency-sensitive, or sparsely-monitored contexts, the gap between compromise and revocation may span hours or days, during which the compromised agent has uninterrupted authority.
\end{enumerate}

\subsection{The Vouch Protocol Approach}

The Vouch Protocol addresses these failures with three composing layers:

\begin{enumerate}[label=\textbf{\arabic*.},leftmargin=*]
\item A \textbf{credential layer} in which every agent action is issued as a W3C Verifiable Credential signed by the agent's Decentralized Identifier (DID) via a Data Integrity proof. The credential's \verb|credentialSubject.intent| field carries the action verb, the target identifier, and the resource URI bound together. The credential's \verb|validUntil| is short (commonly 5 minutes). The proof is \texttt{eddsa-jcs-2022} over the JCS-canonicalized credential bytes, optionally paired with a second \texttt{mldsa44-jcs-2026} proof on the same credential for the dual-proof post-quantum profile.

\item A \textbf{delegation layer} in which credentials chain together, each link cryptographically attesting to a principal-to-sub-principal authorization. Resource scope must narrow at each link; chain depth is bounded; the root principal is verifiable. Multi-agent systems gain end-to-end attribution: ``the patient authorized the clinician, who authorized their EHR agent, who authorized the prior-auth sub-agent, who submitted this credential'' is a single auditable chain.

\item A \textbf{state verifiability layer} in which long-running agents renew their credentials on a heartbeat schedule. Trust decays exponentially over time; behavioral attestation digests are computed per interval; canary commit/reveal chains detect silent failures; an $M$-of-$N$ validator quorum distributes trust evaluation across role-specialized policy / behavior / budget validators. An agent that cannot heartbeat, because it crashed, is compromised, or has drifted out of policy, loses authority by entropy decay; a missed heartbeat is cryptographically detectable.
\end{enumerate}

Together, the three layers move from ``the agent had a key'' to ``the agent issued a signed credential at time $T$, declaring action $A$ on resource $R$, with the authorization of principals $P_1 \to P_2 \to P_n$, verifiable against the agent's published DID Document and the issuer's revocation registry, with a freshness window of 300 seconds and a running behavioral attestation showing the agent is operating within its declared scope.''

\subsection{Design Goals}

Six design goals shaped the protocol:

\begin{itemize}[leftmargin=*]
\item \textbf{Standards alignment over invention.} Where existing open standards suffice (W3C VC 2.0, Data Integrity, DIDs, BitstringStatusList, RFC 8785 JCS, NIST FIPS 204), the protocol composes with them and does not introduce new primitives.
\item \textbf{LLM-isolated key material.} The agent's private signing key must never appear in the LLM's context window. The Identity Sidecar pattern isolates the key in a separate process that the LLM cannot reach.
\item \textbf{Byte-identical canonicalization across languages.} Credentials signed by the Python SDK must verify against credentials signed by the TypeScript or Go SDK with the same input. JCS (RFC 8785) is the canonical form.
\item \textbf{Cryptographic agility.} The protocol must accommodate the post-quantum transition without changing the canonical payload format. The dual-proof profile attaches two independent Data Integrity proofs on the same JCS-canonicalized bytes; verifiers select the trust level they require.
\item \textbf{Continuous verifiability, not point-in-time authentication.} Trust must be renewed under observation; an agent that goes silent must lose authority automatically.
\item \textbf{Open prior art.} Every novel design pattern is published as a defensive disclosure (Apache 2.0 code, CC0 disclosure) to keep the design space open.
\end{itemize}

\subsection{Contribution Summary}

This paper contributes:

\begin{enumerate}[label=\arabic*.,leftmargin=*]
\item The Vouch Credential format and its Data Integrity proof construction (\texttt{eddsa-jcs-2022}).
\item The Identity Sidecar pattern with intent allow-list enforcement at the sidecar layer, defended against prompt injection by construction.
\item The delegation chain mechanism with resource narrowing, depth-limit, and trusted-principal anchoring.
\item The BitstringStatusList integration for per-credential revocation, plus DID-level revocation for issuer-key compromise.
\item The dual-proof post-quantum profile, in which a credential carries two independent Data Integrity proofs (\texttt{eddsa-jcs-2022} and \texttt{mldsa44-jcs-2026}) on the same JCS-canonicalized bytes.
\item The State Verifiability runtime: trust entropy decay, behavioral attestation digests, canary commit/reveal chains, $M$-of-$N$ validator quorum.
\item Cross-language test vectors and three reference implementations (Python, TypeScript, Go).
\end{enumerate}

\subsection{Paper Organization}

Section~\ref{sec:related} reviews related work and clarifies Vouch's position in the agent-identity landscape. Section~\ref{sec:credential} defines the credential format. Section~\ref{sec:sidecar} presents the Identity Sidecar architecture. Section~\ref{sec:delegation} develops the delegation chain construction and its verification. Section~\ref{sec:revocation} details the revocation mechanisms. Section~\ref{sec:hybrid} presents the dual-proof post-quantum profile. Section~\ref{sec:state} introduces the State Verifiability layer. Section~\ref{sec:security} analyses the protocol's security properties against an adversarial threat model. Section~\ref{sec:interop} reports cross-language interoperability results. Section~\ref{sec:discussion} discusses limitations and future work. Section~\ref{sec:conclusion} concludes.

% ====================================================================
\section{Related Work and Positioning}
\label{sec:related}

\subsection{Workload Identity Systems}

SPIFFE and SPIRE~\cite{strong2018spiffe} define workload identity for service-to-service communication in cloud environments. SPIFFE issues short-lived X.509 SVIDs or JWT-SVIDs to workloads via attestation against the host platform. SPIFFE establishes \emph{that} a workload is who it claims to be (via attestation against the runtime), but does not bind the workload's actions to the identity at the per-action level. A SPIFFE-authenticated workload making an HTTP call to a downstream service authenticates the call via mutual TLS; the call's intent (target resource, requested action, authorization chain) is not part of the authenticated payload.

Vouch composes with SPIFFE: a SPIFFE-attested workload may also be a Vouch issuer, holding a Vouch DID whose verification methods are bootstrapped from the SVID. The Vouch credential then provides the per-action binding that SPIFFE does not.

\subsection{OAuth 2.0 / OAuth 2.1 and Token-Based Access}

OAuth 2.x~\cite{rfc6749,rfc9700} is the dominant access-grant protocol for HTTP APIs. Access tokens are bearer credentials authorizing the client to call defined scopes on the resource server. OAuth makes no claim about which specific agent within a client used the token, what the agent intended, or whether the agent's runtime state remains within policy.

Vouch is complementary to OAuth at the application layer. A Vouch credential may accompany an OAuth-authorized HTTP request, providing the agent-identity, intent, and delegation that OAuth's bearer model omits. Resource servers may verify both (OAuth for coarse access authorization and Vouch for action-specific identity and intent), or may use Vouch alone for agent-to-agent communication where OAuth's session model does not fit.

\subsection{Decentralized Identifiers and Verifiable Credentials}

The W3C Decentralized Identifiers (DID) Core~\cite{w3cdidcore} and Verifiable Credentials Data Model 2.0~\cite{w3cvcdm2} specifications define the underlying identity and credential primitives Vouch uses. Vouch issues VCs as defined by VC-DM-2.0, with DIDs as the issuer identifier. Vouch does not introduce a new identity primitive; it specifies a credential subtype (\texttt{VouchCredential}) and a per-action issuance pattern.

Among the many credential subtypes the VC ecosystem has produced (student transcripts, professional certifications, vaccination records, employment letters), Vouch occupies a distinct niche: per-action, short-validity, agent-issued credentials describing what the agent is about to do, not what it has been authorized to be.

\subsection{ZCAP-LD and Capability-Based Authorization}

ZCAP-LD~\cite{zcapld} is a JSON-LD capability authorization format that supports attenuation through delegation. ZCAP capabilities are attenuated at each delegation: a parent capability is restricted as it is delegated downward. The result is similar in shape to a Vouch delegation chain.

Vouch's delegation differs in three respects:

\begin{enumerate}[label=\arabic*.,leftmargin=*]
\item \textbf{Canonicalization}: Vouch uses RFC 8785 JCS, not JSON-LD canonicalization. JCS is simpler, has no external context dependencies, and produces byte-identical output across implementations more reliably.
\item \textbf{Resource binding}: Each Vouch delegation link explicitly binds a \verb|resource| URI. The chain validator enforces that child resources be sub-paths of parent resources, preventing capability widening at any link.
\item \textbf{Chain semantics}: Vouch delegations are credentials in their own right (each link is a VC); ZCAP capabilities are stand-alone documents with their own data model.
\end{enumerate}

\subsection{C2PA Content Credentials}

The Coalition for Content Provenance and Authenticity (C2PA) defines content-binding credentials for media assets~\cite{c2pa14}. C2PA Content Credentials embed signed provenance into media files (images, audio, video). Vouch is complementary: a content-creation agent may sign a Vouch credential authorizing the act of generating an image, then bind the image's C2PA manifest to the same agent identity.

\subsection{Post-Quantum Signature Migration Frameworks}

NIST's selection of ML-DSA-44 (FIPS 204) and SLH-DSA (FIPS 205) as standardized post-quantum digital signature algorithms~\cite{nistfips204} creates the migration pressure Vouch addresses. The IETF's \emph{draft-ietf-jose-pq-composite-sigs} defines composite signatures for JOSE; the W3C Data Integrity community has parallel work in progress, including Digital Bazaar's \texttt{mldsa44-rdfc-2024-cryptosuite} family~\cite{db-mldsa44} and an upcoming JCS variant. Vouch's \emph{dual-proof profile} is a Data-Integrity-native realization that avoids inventing a composite cryptosuite identifier: a single Vouch credential carries two independent Data Integrity proofs in its \verb|proof| array, one \texttt{eddsa-jcs-2022} proof and one \texttt{mldsa44-jcs-2026} proof, both computed over the same JCS canonical credential bytes. Verifiers in classical-only mode validate the \texttt{eddsa-jcs-2022} proof and ignore the rest; verifiers in PQ-aware mode validate the \texttt{mldsa44-jcs-2026} proof; verifiers in both-required mode iterate the proof array and accept only when every proof verifies.

% ====================================================================
\section{The Vouch Credential}
\label{sec:credential}

\subsection{Format}

A Vouch Credential is a W3C Verifiable Credential 2.0 with a \texttt{VouchCredential} type tag, a \verb|credentialSubject| carrying the agent's intent, and a Data Integrity proof over the JCS-canonicalized credential bytes. The format is defined in Section 5 of the Vouch Protocol specification.

\begin{lstlisting}[language=json]
{
  "@context": [
    "https://www.w3.org/ns/credentials/v2",
    "https://vouch-protocol.com/contexts/v1"
  ],
  "id": "urn:uuid:b8b5c7bd-8271-4805-8973-70968c4dd46f",
  "type": ["VerifiableCredential", "VouchCredential"],
  "issuer": "did:web:agent.example.com",
  "validFrom": "2026-05-13T05:41:10Z",
  "validUntil": "2026-05-13T05:46:10Z",
  "credentialSubject": {
    "id": "did:web:agent.example.com",
    "vouchVersion": "1.0",
    "intent": {
      "action": "submit_claim",
      "target": "claim:HC-001",
      "resource": "https://insurance.example.com/claims/HC-001"
    }
  },
  "proof": {
    "type": "DataIntegrityProof",
    "cryptosuite": "eddsa-jcs-2022",
    "created": "2026-05-13T05:41:10Z",
    "verificationMethod": "did:web:agent.example.com#key-1",
    "proofPurpose": "assertionMethod",
    "proofValue": "z77CAhFw1rKB1wLQ541oZ55WD1rVcmkiHFnF8EcVs2A4..."
  }
}
\end{lstlisting}

The required fields beyond the VC 2.0 baseline:

\begin{itemize}[leftmargin=*]
\item \verb|type| MUST include \verb|VouchCredential|.
\item \verb|credentialSubject.vouchVersion| MUST be present and equal to the protocol version this credential conforms to.
\item \verb|credentialSubject.intent.action|, \verb|intent.target|, \verb|intent.resource| MUST all be non-empty strings. Empty strings are treated as missing.
\item \verb|validFrom| MUST be in RFC 3339 form with \verb|Z| UTC indicator (not \verb|+00:00|).
\item \verb|validUntil| SHOULD be set to a short window (default 300 seconds from \verb|validFrom|).
\end{itemize}

\subsection{The eddsa-jcs-2022 Cryptosuite}

The default cryptosuite is \texttt{eddsa-jcs-2022} (W3C Data Integrity registered identifier). The signing procedure:

\begin{enumerate}[label=\arabic*.,leftmargin=*]
\item Construct the credential dictionary with all required fields, including the \verb|proof| object except \verb|proofValue|.
\item Canonicalize the entire credential using RFC 8785 JCS. JCS specifies a strict deterministic ordering, escaping, and number formatting; the output is byte-identical across conforming implementations.
\item Compute the Ed25519 signature over the canonical bytes using the issuer's private key.
\item Encode the signature in multibase base58btc form (the \texttt{z} prefix), producing approximately 88 characters.
\item Set \verb|proof.proofValue| to the multibase-encoded signature.
\end{enumerate}

Verification reverses the process:

\begin{enumerate}[label=\arabic*.,leftmargin=*]
\item Extract \verb|proof.proofValue|, decode from multibase base58btc.
\item Reconstruct the credential dictionary excluding \verb|proof.proofValue| (other proof fields are part of the signed payload).
\item JCS-canonicalize.
\item Resolve the issuer's DID Document, locate the verification method indicated by \verb|proof.verificationMethod|, extract the Ed25519 public key (encoded as Multikey in the DID Document).
\item Verify the Ed25519 signature.
\end{enumerate}

\subsection{JCS Determinism Across Languages}

JSON Canonicalization Scheme (RFC 8785) fixes the serialization that Ed25519 signs. The protocol's three reference implementations (Python, TypeScript, Go) produce byte-identical JCS output given byte-identical inputs. Cross-language test vectors verify the property on every release.

The byte-identity property is foundational: without it, a credential signed by Python and a credential with the same logical content signed by TypeScript would have different signatures, breaking interoperability and federation across implementations.

\subsection{Intent Replay Resistance}

Each credential carries a unique \verb|id| (typically a UUID URN). Verifiers maintain a nonce store keyed on \verb|id|. A credential whose \verb|id| has previously been seen is rejected (\verb|nonce_replay|). The nonce store's TTL must be at least the longest plausible credential \verb|validUntil| to prevent a credential from being replayed after the nonce store has forgotten it.

Replay resistance is also enforced at the intent level: \verb|intent.target| and \verb|intent.resource| are part of the signed payload. A credential authorizing \verb|submit_claim| on \verb|claim:HC-001| cannot be replayed against \verb|claim:HC-002| without re-signing; the JCS canonical bytes differ.

% ====================================================================
\section{The Identity Sidecar Pattern}
\label{sec:sidecar}

\subsection{Motivation}

Modern AI agents are typically built around a Large Language Model. The LLM is stochastic, non-deterministic, and vulnerable to prompt injection. Exposing the agent's private signing key to the LLM's context window creates three failure modes:

\begin{itemize}[leftmargin=*]
\item \textbf{Key leak via prompt injection.} Adversarial content in retrieved documents, tool outputs, or user input can instruct the LLM to print the key as part of its output.
\item \textbf{Unauthorized signing via jailbreak.} A jailbroken LLM may use the key to sign arbitrary intents, including intents not authorized by the principal.
\item \textbf{Persistence in training data and logs.} Key material that passed through the LLM may end up in logged conversation transcripts, fine-tuning corpora, or vendor-side analytics.
\end{itemize}

Direct key exposure to the LLM is incompatible with the protocol's security model.

\subsection{Architecture}

The Identity Sidecar pattern separates the agent into two processes:

\begin{enumerate}[label=\arabic*.,leftmargin=*]
\item \textbf{The Brain (stochastic).} The LLM. Holds zero cryptographic secrets. Reasons about which action to take and proposes intents to the Sidecar.
\item \textbf{The Passport (deterministic).} A small, auditable process holding the private signing keys in its address space. Receives intent proposals from the Brain, applies a deterministic policy check, and issues credentials only when policy passes.
\end{enumerate}

The Brain and the Passport communicate over a local IPC channel: typically localhost HTTP, a UNIX domain socket, or an MCP transport. The IPC boundary is the trust boundary.

\subsection{Just-In-Time Signing Flow}

A signing operation proceeds:

\begin{enumerate}[label=\arabic*.,leftmargin=*]
\item The Brain decides an action is appropriate. It constructs a proposed intent.
\item The Brain submits the intent to the Sidecar over IPC.
\item The Sidecar applies its policy: structural validation (required fields present), allow-list check (action type is in the deployed allow-list), and rate-limit check.
\item If policy passes, the Sidecar constructs the full credential, signs it with the held private key, and returns the credential to the Brain.
\item If policy fails, the Sidecar returns a structured error. No credential is produced.
\end{enumerate}

The Sidecar's policy check is the \emph{last line of defence} against a compromised Brain. A prompt-injected LLM may propose any intent it likes; the Sidecar refuses to sign anything outside the allow-list. The capability bound on a compromised Brain is therefore structural, not behavioural.

\subsection{Reference Implementations and Tier Hierarchy}

Three reference Sidecar implementations accompany the protocol:

\begin{itemize}[leftmargin=*]
\item \textbf{Go} (\verb|go-sidecar/|): production-tier, KMS/HSM-backed, hybrid PQ supported, sensitive-mode JWE wrapping, multi-tenant.
\item \textbf{Python} (\verb|vouch.sidecar.*|): lightweight tier for self-hosted Python stacks. File or environment keys. Intentionally omits KMS, hybrid PQ, JWE, and multi-tenancy.
\item \textbf{TypeScript} (\verb|packages/sdk-ts/sidecar/|): lightweight tier for Node stacks. Same omissions as Python.
\end{itemize}

The HTTP wire contract is identical across all three. A shared contract test suite verifies that all three implementations accept and reject the same inputs and emit semantically equivalent credentials. The tier hierarchy is informative guidance: deployers requiring KMS, FIPS, hybrid PQ, or multi-tenancy run the Go sidecar; deployers with simpler requirements may run Python or TypeScript.

% ====================================================================
\section{Delegation Chains}
\label{sec:delegation}

\subsection{Motivation}

Multi-agent systems present an attribution problem: an action taken by a leaf agent may have been authorized by a chain of principals: a user authorized an assistant, which authorized a research agent, which authorized a web-scraping sub-agent. When the action requires audit, the leaf-only signing model loses the chain.

Vouch delegations are themselves Verifiable Credentials. Each link in a chain is a credential signed by the delegating principal and identifying the delegate. A leaf-action credential includes the full chain (or a reference to it) so a verifier can reconstruct the authorization path from root principal to leaf agent.

\subsection{Delegation Link Structure}

A delegation link is a Vouch credential whose \verb|credentialSubject| is shaped:

\begin{lstlisting}[language=json]
{
  "credentialSubject": {
    "id": "did:web:research-agent.example.com",
    "delegatedBy": "did:web:assistant.example.com",
    "scope": {
      "actions": ["web_search", "scrape_page"],
      "resource": "https://research-agent.example.com/agents/research-agent/scope/2026"
    },
    "validUntil": "2026-05-14T08:00:00Z"
  }
}
\end{lstlisting}

The link's \verb|issuer| is the delegating principal. The link's \verb|credentialSubject.id| is the delegate. \verb|scope.actions| enumerates which actions are delegated; \verb|scope.resource| is the resource URI under which all delegated actions must fall.

\subsection{Chain Construction and the Resource Narrowing Rule}

A chain is a list of delegation-link credentials, ordered from root principal to leaf agent. At each link, the child's \verb|scope.resource| MUST be a sub-URI of the parent's \verb|scope.resource|. The narrowing rule prevents capability widening at any link. A delegate cannot grant a sub-delegate more authority than it itself holds.

\subsection{Depth Limit}

Chain depth is bounded to \textbf{five links maximum}. This limit prevents pathological chains from inflating verification cost and from obscuring the actual authorization path.

\subsection{Chain Validation}

A verifier validating a leaf credential with an attached chain:

\begin{enumerate}[label=\arabic*.,leftmargin=*]
\item Verifies the leaf credential's signature.
\item Identifies the chain root: the first link's \verb|issuer| MUST be in the verifier's set of trusted principals for the leaf action.
\item For each adjacent pair (parent, child) in the chain, verifies the child's \verb|scope.resource| is a sub-URI of the parent's.
\item Verifies each link's \verb|delegatedBy| matches the previous link's \verb|credentialSubject.id|.
\item Verifies each link's signature.
\item Verifies the leaf credential's \verb|intent.action| is in the deepest link's \verb|scope.actions|.
\item Verifies the leaf credential's \verb|intent.resource| is a sub-URI of the deepest link's \verb|scope.resource|.
\end{enumerate}

Any failure produces a structured rejection reason: \verb|untrusted_principal|, \verb|chain_depth_exceeded|, \verb|resource_not_narrowed|, \verb|link_signature_invalid|, \verb|parent_proof_mismatch|, or \verb|action_not_in_scope|.

% ====================================================================
\section{Revocation}
\label{sec:revocation}

\subsection{Two Mechanisms}

Vouch supports two complementary revocation mechanisms:

\begin{itemize}[leftmargin=*]
\item \textbf{DID-level revocation}: revoke an entire DID, invalidating all credentials it has issued or ever will issue.
\item \textbf{Credential-level revocation}: revoke a specific credential without affecting other credentials from the same issuer.
\end{itemize}

Most production deployments run both. DID-level revocation is appropriate for blanket kill switches (key compromise, agent decommissioning); credential-level revocation is appropriate for surgical retraction (regulatory hold on a specific transaction, suspension pending review).

\subsection{DID-Level Revocation Registry}

The protocol defines a revocation registry interface (\verb|RevocationStoreInterface|). A revoked DID is recorded as a \verb|RevocationRecord| with \verb|did|, \verb|revoked_at|, \verb|reason|, and \verb|revoked_by|. Reference implementations include in-memory, file, and Redis backends.

Verifiers consult the registry on every credential verification. If the issuing DID is in the registry, the credential is rejected with \verb|issuer_revoked|. The registry's cache TTL is operator-configurable; a typical default is 60 seconds.

\subsection{BitstringStatusList for Per-Credential Revocation}

For per-credential revocation, Vouch uses W3C BitstringStatusList~\cite{w3cbsl}. The issuer maintains a published \texttt{BitstringStatusListCredential} whose \verb|credentialSubject.encodedList| is a gzip-compressed, base64url-encoded bitstring. Each Vouch credential includes a \verb|credentialStatus| property pointing at its index. A verifier fetches the status list, decompresses, reads the indicated bit; if the bit is set, the credential is revoked.

The protocol's \verb|StatusListFetcher| caches the status list with a configurable TTL and supports conditional GETs (\verb|If-None-Match|, \verb|If-Modified-Since|) for efficient re-validation.

% ====================================================================
\section{Dual-Proof Post-Quantum Profile}
\label{sec:hybrid}

\subsection{Motivation}

NIST's selection of ML-DSA-44 (FIPS 204) as a standardized post-quantum digital signature~\cite{nistfips204} places a known cryptographic transition on the protocol's horizon. Ed25519 will not survive a sufficiently large fault-tolerant quantum computer. The transition window is uncertain but real, and credentials with multi-year audit retention requirements (regulated healthcare, financial settlement) cannot afford to be signed with an algorithm that will be broken before the retention window expires.

The protocol's response is the \emph{dual-proof profile}: a credential carries two independent W3C Data Integrity proofs on the same JCS-canonicalized credential bytes, one \texttt{eddsa-jcs-2022} proof (classical Ed25519) and one \texttt{mldsa44-jcs-2026} proof (ML-DSA-44). The Data Integrity specification already defines the \verb|proof| field as an array, so the dual-proof construction is a natural use of existing primitives and requires no Vouch-specific composite cryptosuite identifier.

\subsection{Construction}

The credential's \verb|proof| field is an array of two Data Integrity proof objects:

\begin{verbatim}
"proof": [
  {
    "type": "DataIntegrityProof",
    "cryptosuite": "eddsa-jcs-2022",
    "verificationMethod": "...#key-ed25519",
    "proofPurpose": "assertionMethod",
    "proofValue": "z<base58btc(ed25519_sig)>"
  },
  {
    "type": "DataIntegrityProof",
    "cryptosuite": "mldsa44-jcs-2026",
    "verificationMethod": "...#key-mldsa44",
    "proofPurpose": "assertionMethod",
    "proofValue": "z<base58btc(mldsa44_sig)>"
  }
]
\end{verbatim}

Both proofs are computed over the same JCS-canonical credential bytes (the credential minus the \verb|proofValue| field of each proof, JCS-canonicalized). The same-canonical-bytes property is preserved by construction: both cryptosuites use JCS canonicalization on the same credential body. Each proof's \verb|proofValue| is independently multibase base58btc-encoded.

The Ed25519 signature is exactly 64 bytes; the ML-DSA-44 signature is exactly 2,420 bytes. Each appears in its own proof object's \verb|proofValue| rather than being concatenated into a single field.

\subsection{Verifier Modes}

A verifier iterates the credential's \verb|proof| array and applies one of three local policies:

\begin{enumerate}[label=\arabic*.,leftmargin=*]
\item \textbf{Classical-only}: validate the \texttt{eddsa-jcs-2022} proof; ignore the rest. Useful when ML-DSA-44 libraries are unavailable or computationally expensive.
\item \textbf{Post-quantum-only}: validate the \texttt{mldsa44-jcs-2026} proof. Useful when classical signatures are no longer trusted under the verifier's compliance regime.
\item \textbf{Both-required}: validate every proof in the array; reject if any one fails. The cautious mode, recommended for long-retention credentials.
\end{enumerate}

The verifier's mode is policy, not part of the credential. The same credential is verifiable under any verifier mode without re-signing.

\subsection{DID Document Representation}

A signing DID supporting the dual-proof profile publishes two verification methods in its DID Document: one Ed25519 Multikey, one ML-DSA-44 Multikey. The Multikey multicodec prefix distinguishes the two algorithms, and each proof's \verb|verificationMethod| field points at its own key.

\subsection{Migration Path}

A deployment migrates from classical to dual-proof in three steps:

\begin{enumerate}[label=\arabic*.,leftmargin=*]
\item \textbf{Adopt dual-proof signing}: the signer issues credentials with two proofs in the array. Existing classical-only verifiers continue working (they see a familiar \texttt{eddsa-jcs-2022} proof and ignore the second one they don't recognize).
\item \textbf{Update verifiers}: deploy a verifier release that supports ML-DSA-44. Configure the verifier's mode (classical-only, PQ-only, or both-required) based on the deployment's risk tolerance.
\item \textbf{Retire classical signing}: once all verifiers are PQ-aware, switch the signer to issue only the \texttt{mldsa44-jcs-2026} proof. The dual-proof profile is the migration vehicle, not the destination.
\end{enumerate}

\subsection{Performance and Size}

Per-credential signing cost (Apple M3 Max, reference Python implementation):

\begin{itemize}[leftmargin=*]
\item Ed25519 alone: $\sim$50 $\mu$s
\item ML-DSA-44 alone: $\sim$3 ms
\item Dual-proof (both): $\sim$3 ms (dominated by ML-DSA-44)
\end{itemize}

Per-credential credential size (multibase-encoded proofs):

\begin{itemize}[leftmargin=*]
\item Classical only: $\sim$700 bytes total
\item Dual-proof: $\sim$3{,}200 bytes total
\end{itemize}

The size delta matters for HTTP-header-conveyed credentials and for credentials embedded in QR codes. For application-body conveyance the overhead is negligible.

\subsection{Relationship to Concurrent Work}

The dual-proof construction converges with Digital Bazaar's \texttt{mldsa44-rdfc-2024-cryptosuite} family~\cite{db-mldsa44}, which defines an ML-DSA-44 Data Integrity cryptosuite using RDFC canonicalization, and with that family's forthcoming JCS variant. The cryptosuite identifier \texttt{mldsa44-jcs-2026} used in this paper is provisional and is being aligned with that upstream registration. Earlier Vouch reference implementations (v1.6.x) emit a single composite cryptosuite (\texttt{hybrid-eddsa-mldsa44-jcs-2026}) with a concatenated \verb|proofValue|; this composite identifier is retained as a transitional alias for backward compatibility only, and new implementations emit the dual-proof form described above.

% ====================================================================
\section{The State Verifiability Layer}
\label{sec:state}

\subsection{Motivation: From Point-in-Time to Continuous}

The credential layer addresses ``did the agent claim authorization at time $T$.'' The State Verifiability layer addresses operational questions arising \emph{after} that: is the agent still behaving correctly \emph{now}? Has it drifted out of policy? Has it been substituted? Has it crashed silently and been replayed by an adversary?

A long-running agent without continuous attestation accumulates risk. The longer it runs, the wider the window during which a compromise or substitution may go undetected. Heartbeat-style renewal makes that window bounded.

\subsection{The Heartbeat Protocol}

The Heartbeat Protocol defines a periodic renewal cycle:

\begin{enumerate}[label=\arabic*.,leftmargin=*]
\item The agent's \verb|HeartbeatSession| records actions and behavioural signals during each interval.
\item At the heartbeat boundary (default 60 seconds), the session constructs a \verb|HeartbeatRequest|: it includes a behavioural digest, the running Merkle root of actions performed since the last heartbeat, the canary commit/reveal pair, and an interval index.
\item The agent submits the request to a validator (or quorum of validators).
\item The validator(s) check schema, behavioural digest structure, canary chain integrity, interval-index monotonicity, and trust policy.
\item On success, the validator returns a fresh \verb|SessionVoucher| credential with the agent's renewed trust parameters.
\item On failure (broken canary, stale interval index, behavioural drift), no voucher is issued; the agent's existing voucher expires; the agent loses authority.
\end{enumerate}

The Heartbeat Protocol inverts the traditional PKI trust model from ``trusted until revoked'' to \emph{``untrusted until renewed.''}

\subsection{Trust Entropy Decay}

Each SessionVoucher carries \texttt{initialTrust} and \texttt{decayLambda}. The effective trust at time $t$ after the voucher's \texttt{issuedAt}:

\[
\textsf{trust}(t) = \textsf{initialTrust} \cdot e^{-\lambda (t - t_0)}
\]

Sensitive actions are gated by current trust. The protocol defines three reference thresholds:

\begin{itemize}[leftmargin=*]
\item \verb|TRUST_THRESHOLD_HIGH_STAKES = 0.9|: financial transfers, code deployment, regulated submissions.
\item \verb|TRUST_THRESHOLD_MEDIUM_STAKES = 0.75|: PHI read, customer-data access.
\item \verb|TRUST_THRESHOLD_LOW_STAKES = 0.5|: status queries, idle activity.
\end{itemize}

Heartbeat interval should be set to less than half the half-life so renewal stays ahead of decay. For $\lambda = 0.01$, half-life is approximately 69 seconds; a 30-second heartbeat is appropriate.

\subsection{Behavioural Attestation}

Per-interval, the \verb|BehavioralCollector| records:

\begin{itemize}[leftmargin=*]
\item API calls made by the agent (URL, token count, optional intent-drift score)
\item Resources accessed
\item Any policy-defined custom signals
\end{itemize}

At each heartbeat, the collector produces a \verb|behavioralDigest|. Three reference intent-drift scorers ship with the SDK: arithmetic mean, maximum-of-samples, and exponential-weighted moving average. The drift score quantifies the divergence of the agent's recent activity from a baseline; a sharp rise across consecutive intervals signals a potential prompt-injection or policy-drift event.

\subsection{Canary Commitments}

The canary commit/reveal chain provides cryptographic detection of silent agent failure. Each heartbeat commits to a fresh 256-bit secret hash; the next heartbeat reveals the prior secret. A missed heartbeat means the prior secret is never revealed; no subsequent heartbeat can resume the chain because the validator expects the missing reveal.

The validator state per agent is one string (the last commitment). The state is cheap to persist and survives validator restarts.

\subsection{Action Merkle Roots}

A \verb|HeartbeatRequest| includes the Merkle root of actions performed since the last heartbeat, computed over RFC 6962-domain-separated Merkle trees~\cite{rfc6962}. Each action is hashed with \texttt{0x00} prefix (leaf); internal nodes are hashed with \texttt{0x01} prefix. This prevents the classical second-preimage attack on Merkle trees.

\subsection{Validator Quorum}

Single validators are single points of failure. The \verb|HeartbeatQuorum| distributes trust evaluation across multiple validators with different responsibilities:

\begin{lstlisting}[language=Python]
quorum = HeartbeatQuorum(
    validators=[
        QuorumValidator(validator=policy_validator,     role=ROLE_POLICY),
        QuorumValidator(validator=behavioral_validator, role=ROLE_BEHAVIORAL),
        QuorumValidator(validator=budget_validator,     role=ROLE_BUDGET),
    ],
    threshold=2,
)
\end{lstlisting}

Roles are reference tags; deployments choose their own role taxonomy. The quorum threshold is $M$-of-$N$. Trust parameter aggregation across approving validators is configurable; the defaults bias toward caution.

% ====================================================================
\section{Security Analysis}
\label{sec:security}

\subsection{Threat Model}

The protocol assumes:

\begin{itemize}[leftmargin=*]
\item The cryptographic primitives (Ed25519, ML-DSA-44, SHA-256) are sound under their respective security assumptions.
\item The verifier can resolve the issuer's DID Document via the agent's chosen DID method (\texttt{did:web} via HTTPS, \texttt{did:key} via decoding).
\item The verifier has access to the revocation registry and to BitstringStatusList endpoints.
\item The Identity Sidecar's process boundary is enforced by the host operating system: a malicious LLM process cannot read the Sidecar's memory.
\end{itemize}

The protocol does \emph{not} assume:

\begin{itemize}[leftmargin=*]
\item The LLM is trustworthy. Prompt injection, jailbreak, and model substitution attacks are within the threat model.
\item The transport layer is confidential. Credentials are designed to be exposed in transit (the signature establishes authenticity even when an adversary observes the credential).
\item The agent's host platform is uncompromised. Root-level compromise of the host defeats any application-level defence, including this one.
\end{itemize}

\subsection{Defended Threats}

\noindent\textbf{Bearer-credential theft.} Vouch credentials are bound to the issuer's DID by signature; an adversary who copies a credential cannot reuse it to issue new credentials. The original credential is bound to its \verb|intent.target| and \verb|intent.resource|; it cannot be replayed against a different resource.

\noindent\textbf{Prompt injection causing key exfiltration.} The Identity Sidecar makes this attack class structurally impossible: the LLM process has no access to the private key.

\noindent\textbf{Prompt injection causing unauthorized credential issuance.} The Sidecar's intent allow-list makes this attack class capability-bounded: the LLM may propose any intent it likes, but the Sidecar refuses to sign intents outside the deployment's allow-list.

\noindent\textbf{Replay across resources.} Each credential's \verb|intent.resource| is part of the signed payload. Replaying against a different resource invalidates the signature.

\noindent\textbf{Delegation chain forgery.} Each link is a signed credential; forging a link requires possessing the delegating principal's key. Trusted-principal anchoring ensures the chain root is known to the verifier independently.

\noindent\textbf{Resource widening at delegation.} The chain validator's resource-narrowing rule rejects chains where any link grants access beyond its parent's scope.

\noindent\textbf{Long-running agent silent failure.} The Heartbeat Protocol's canary chain detects missed heartbeats: the next heartbeat cannot produce the expected reveal, the chain breaks, and the validator refuses to renew.

\noindent\textbf{Key compromise.} DID-level revocation invalidates all credentials issued by the compromised key. The auto-rotation pipeline automates rotation when leak detection fires.

\noindent\textbf{Post-quantum cryptanalysis.} The dual-proof profile carries an ML-DSA-44 Data Integrity proof in addition to the Ed25519 proof, allowing the credential to be re-verified against the PQ proof once Ed25519 is broken. Verifiers in both-required mode are immune to single-algorithm cryptanalysis.

\subsection{Residual Threats}

\noindent\textbf{Within-allow-list abuse.} The Sidecar's allow-list bounds capability \emph{type}, not \emph{abuse within type}. Mitigations: pattern strictness, rate limits, downstream verifier policy, and human-in-the-loop confirmation for ambiguous cases.

\noindent\textbf{DID resolution failures.} If the issuer's DID Document is unreachable, the credential cannot be verified. Verifiers MUST fail closed rather than accept unverified credentials.

\noindent\textbf{Trusted-principal anchor compromise.} If a trusted-principal DID is compromised, the adversary can construct delegation chains rooted at that DID that authorize arbitrary leaf actions within the principal's scope. Defence: rotate trusted-principal keys regularly; require multi-key signing at the root for high-stakes deployments.

% ====================================================================
\section{Cross-Language Interoperability}
\label{sec:interop}

\subsection{Test Vectors}

The protocol ships cross-language test vectors verifying that Python, TypeScript, and Go implementations produce byte-identical credentials given identical inputs:

\begin{itemize}[leftmargin=*]
\item \verb|test-vectors/jcs/|: JSON Canonicalization Scheme reference outputs.
\item \verb|test-vectors/eddsa-jcs-2022/|: full credential signing test vectors for the classical cryptosuite.
\item \verb|test-vectors/hybrid-eddsa-mldsa44/|: full credential signing test vectors for the dual-proof post-quantum profile (and, for v1.6.x interop, the transitional composite cryptosuite).
\item \verb|test-vectors/bitstring-status-list/|: BitstringStatusList encoding test vectors.
\item \verb|test-vectors/sidecar-contract/|: HTTP wire-contract test suite for the three sidecar implementations.
\end{itemize}

\subsection{Reference Implementation Properties}

\begin{table}[h]
\centering
\small
\begin{tabular}{lllrrl}
\toprule
Implementation & Language & LOC & Tests & Cryptosuites & Sidecar tier \\
\midrule
\verb|vouch/| & Python 3.10+ & 4{,}500 & 350 & classical + hybrid & Lightweight + Dev \\
\verb|packages/sdk-ts/| & TypeScript 5+ & 3{,}200 & 120 & classical + hybrid & Lightweight \\
\verb|go-sidecar/| & Go 1.21+ & 2{,}800 & 75 & classical + hybrid & Production \\
\bottomrule
\end{tabular}
\caption{Reference implementations of the Vouch Protocol.}
\end{table}

All three pass the cross-language test vectors. CI runs the vectors on every commit, and rejects merges that introduce divergence.

\subsection{Known Implementation Divergences}

The only documented divergence is in BitstringStatusList byte-encoding: Python and TypeScript use \texttt{zlib}'s DEFLATE encoder, which produces byte-identical output for identical bitstrings; Go's \texttt{compress/flate} produces a valid but different DEFLATE stream. The spec requires equivalence of the \emph{decompressed} bitstring; all three implementations agree on the decompressed bitstring and therefore on every credential's revocation status. Verifiers and issuers interoperate across all three.

% ====================================================================
\section{Discussion}
\label{sec:discussion}

\subsection{Limitations}

\noindent\textbf{State Verifiability is Python-only at the runtime level.} The data formats (SessionVoucher, behavioural digest, canary commitment, heartbeat request) are cross-language and verifiable in any of the three implementations. The runtime orchestration (\verb|HeartbeatSession|, \verb|HeartbeatScheduler|, \verb|HeartbeatValidator|, \verb|HeartbeatQuorum|) is implemented in Python only. TypeScript and Go ports are work in progress at the time of publication.

\noindent\textbf{Trusted-principal anchoring is deployment-specific.} The protocol does not standardize how a verifier discovers its trusted principals. Each deployment configures this independently.

\noindent\textbf{Confidentiality of sensitive intents is not addressed in v0.1-draft.} Vouch credentials are non-confidential. For deployments handling regulated content, the next minor revision will add an optional confidentiality profile using post-quantum key encapsulation.

\noindent\textbf{The Identity Sidecar's allow-list is static configuration.} Dynamic capability adjustment is achievable through the protocol's normal credential validity windows but requires care to compose with the Sidecar's static allow-list.

\subsection{Adoption Path}

The protocol has been validated against three reference implementations and a cross-language test-vector battery. Adoption requirements for a new deployment:

\begin{enumerate}[label=\arabic*.,leftmargin=*]
\item Generate an issuer DID and publish its DID Document.
\item Choose a Sidecar tier (Go for production, Python or TypeScript for lighter deployments).
\item Configure the Sidecar's allow-list with the deployment's action vocabulary.
\item Wire the agent's tool-call layer to call the Sidecar's \verb|/sign| endpoint instead of directly calling external APIs.
\item Deploy a verifier (either as a library at the API boundary, or as a separate service) for the deployment's external services.
\item For long-running agents, deploy at least one Heartbeat validator (or a quorum of three for regulated deployments).
\end{enumerate}

End-to-end onboarding effort is on the order of one to two engineering days for the credential layer alone, and another two to four days for state verifiability and quorum.

\subsection{Future Work}

Four directions are actively under development:

\begin{enumerate}[label=\arabic*.,leftmargin=*]
\item \textbf{TypeScript and Go ports of the State Verifiability runtime.}
\item \textbf{Confidentiality profile for sensitive intents} using ML-KEM-768 key encapsulation.
\item \textbf{Hosted continuous leak monitor} that closes the loop from automated detection to DID rotation with dual-signed migration credentials and verifier broadcast.
\item \textbf{AI-assisted developer tooling} distributed as Claude Skills / Custom GPTs / Gemini Gems, allowing developers to receive Vouch integration guidance through their own AI tool subscriptions with zero protocol-vendor inference cost.
\end{enumerate}

The protocol's State Verifiability runtime, the dual-proof post-quantum profile, and the migration story from classical to dual-proof to pure-PQ each merit follow-up papers; outlines are in preparation.

% ====================================================================
\section{Conclusion}
\label{sec:conclusion}

We have presented the Vouch Protocol, an open standard for cryptographic identity and continuous state verifiability of autonomous AI agents. The protocol composes Verifiable Credentials, Decentralized Identifiers, Data Integrity proofs, and post-quantum cryptosuites with two novel layers: an Identity Sidecar pattern that isolates the agent's key from the LLM and bounds the agent's capabilities by an enforced allow-list, and a State Verifiability layer that renews trust continuously via heartbeat, decays trust exponentially in the interval, attests behaviour cryptographically, and detects silent failure through canary commit/reveal chains.

Three reference implementations (Python, TypeScript, Go) interoperate at the credential-byte level. The protocol is open under Apache 2.0; sixty defensive disclosures keep the design space open under CC0.

The accountability gap in agent identity is solvable today. The cryptographic primitives are sound, the standards exist, and the implementation effort is bounded. Vouch is one realization of a path the open agent-identity ecosystem will increasingly need: from ``the bearer has access'' to ``this specific agent issued this specific intent under this specific authorization chain, verifiable in seconds, with cryptography rather than logbook evidence.''

% ====================================================================
\section*{Acknowledgements}

The author thanks Manu Sporny and the W3C Verifiable Credentials Working Group for foundational work on Verifiable Credentials, Data Integrity, and the eddsa-jcs-2022 cryptosuite; the W3C Credentials Community Group for hosting and reviewing the incubation; the IETF JOSE working group for ongoing work on PQ/T composite signatures; the C2PA technical committee for content-provenance complementarity; and the open-source community of contributors and reviewers whose feedback shaped many of the choices documented here.

% ====================================================================
\bibliographystyle{plain}
\bibliography{references}

% ====================================================================
\appendix

\section{Reference Implementations}

\begin{itemize}[leftmargin=*]
\item \textbf{Python}: \verb|pip install vouch-protocol|. Source: \url{https://github.com/vouch-protocol/vouch/tree/main/vouch/}
\item \textbf{TypeScript}: \verb|npm install @vouch-protocol/sdk|. Source: \url{https://github.com/vouch-protocol/vouch/tree/main/packages/sdk-ts/}
\item \textbf{Go}: install with \verb|go install| from \url{github.com/vouch-protocol/vouch/go-sidecar/cmd/vouch-sidecar}. Source: \url{https://github.com/vouch-protocol/vouch/tree/main/go-sidecar/}
\end{itemize}

\section{Test Vectors}

All test vectors are at \url{https://github.com/vouch-protocol/vouch/tree/main/test-vectors/}.

\section{Specification Document}

The full Vouch Protocol Specification, currently Spec v0.1-draft in W3C Credentials Community Group incubation, is at \url{https://github.com/vouch-protocol/vouch/blob/main/docs/specs/w3c-cg-report.md}.

\end{document}
